\documentclass[10pt,a4paper]{amsart}
\usepackage[margin=19mm]{geometry}
\usepackage{amsmath,amssymb,mathtools}
\usepackage[scr=rsfs,cal=euler]{mathalfa}
\usepackage{xfrac}
\usepackage[unicode,hidelinks,bookmarks=false]{hyperref}
\newcommand{\ie}{i.e.\ }
\newcommand{\cf}{cf.\ }
\newcommand{\resp}{respectively\ }
\newcommand{\Subsection}{\subsection}
\providecommand{\nobreakdash}{\nobreak\hbox{-}}
\setlength{\emergencystretch}{2em}
\multlinegap0pt
\usepackage{breqn}
\usepackage{amssymb}
\newtheorem{proposition}{Proposition}
\newtheorem{lemma}{Lemma}
\newcommand{\A}{\mathcal{A}}
\newcommand{\B}{\mathcal{B}}
\newcommand{\Le}{\mathcal{L}}
\newcommand{\Lep}{\mathscr{L}}
\newcommand{\N}{\mathbb{N}}
\title{The strong topological Rokhlin property and Medvedev degrees of SFTs}
\author{Nicanor Carrasco-Vargas}
\date{Source: arXiv:2601.03501v1 (CC BY 4.0)}
\begin{document}
\maketitle

\noindent\emph{Source and licence.} The strong topological Rokhlin property and Medvedev degrees of SFTs. Version arXiv:2601.03501v1 (CC BY 4.0), distributed by arXiv under the Creative Commons Attribution 4.0 International licence, http://creativecommons.org/licenses/by/4.0/, as recorded in the arXiv licence metadata for this exact version and shown on the abstract page https://arxiv.org/abs/2601.03501v1. Full licence notice: \texttt{licenses/arxiv-2601.03501-CC-BY-4.0.txt}.\par\smallskip

\noindent\textit{Editorial scope.} Counted unit: the article's proposition prop:projectively-isolated-decidable-language (source0.tex lines 178--180). The article's complete original proof of it is delivered with it (source0.tex lines 181--208, one \texttt{proof} environment of 28 lines). No mathematics is written, repaired or re-derived: the statement and the proof are the article's own lines, in the article's own notation, and no retained line is substituted or re-flowed.  Retained uncounted as the source-local context the proof needs: lines 209--214.  Nothing was omitted from any retained span, so the package carries no omission record.  No retained line contains a citation, so the wrapper carries no bibliography and no bibliography entry.  No retained line contains a figure environment, an image-inclusion command or drawing code, so no image is delivered and the package declares no asset dependency.  Editorial formatting only, in addition: an AMS \texttt{amsart} preamble that restates the article's own declarations and macros used by the retained lines, namely the environments \texttt{proposition}, \texttt{lemma} on separate counters, together with the standard packages those lines need.  The retained environments print in this document as Proposition 1, Lemma 1 in order of appearance, whereas the article prints the same items with the numbering of its own full document.  The complete native source of the article as distributed in the arXiv e-print for this exact version is bundled unchanged as \texttt{sources/192231/source0.tex}.
\begin{lemma}\label{lemma}
    Let $G$ be a finitely generated group and let $\phi\colon \B^G\to \A^G$ be a morphism. Let $X\subset\A^G$ and  $Y\subset \B^G$ be effective subshifts satisfying $X=\phi(Y)$. Let $p\colon F\to \A$ be a pattern presentation with support $F\Subset S^\ast$, and define
    \[X_p=\{x\in X : \text{$p$ does not appear in $x$}\},\] 
    \[Y_p=\{y \in Y : \phi(y)\in X_p\}.\]
    Then $\Lep^c(Y_p)$ is recursively enumerable uniformly on $p$. That is, $\{(q,p) : q\in \Lep^c(Y_p)\}$ is recursively enumerable.  
\end{lemma}

\begin{proposition}\label{prop:projectively-isolated-decidable-language}
A projectively isolated subshift on a recursively presented group has decidable language. 
\end{proposition}

\begin{proof}
    Let $X\subset \A^G$ be a projectively isolated subshift. Our goal is showing that $\Lep(X)$ is decidable. Since a projectively isolated subshift is sofic and sofic subshifts are effective, it follows that $\Lep^c(X)$ is recursively enumerable. We now prove that $\Lep(X)$ is also recursively enumerable. 
    %Thus we exhibit an algorithm which on input a pattern presentation $p\colon F\to \A$, $F\Subset S^\ast$, halts if and only if $p\in\Lep(X)$. 

    

    Fix a finite alphabet $\B$, an open set $U\subset S(\B^G)$, and a morphism $\phi\colon \B^G\to \A^G$ as in the definition of projectively isolated subshift associated to $X$. Since $U$ is open and SFTs are dense in $S(\B^G)$, we can pick an SFT $Y\in U$. Since $U$ is open, we can find a radius $\epsilon$ so that the ball in the metric $D$ with radius $\epsilon$ and center $Y$ is contained in $U$. Fix $k\in\N$ so that $2^{-k}<\epsilon$. It follows then from the definition of $D$ that every subshift $Z\in S(\B^G)$ with $\Le_{B_k}(Z)=\Le_{B_k}(Y)$ satisfies $Z\in U$. 

    In order to prove that $\Lep(X)$ is recursively enumerable, we consider a pattern presentation $p\colon F\to \A$ with support $F\Subset S^\ast$, and describe an algorithmic procedure that halts if and only if $p\in\Lep(X)$. For this we define two subshifts $X_p$ and $Y_p$ which depend on $p$.
    \[X_p=\{x\in X : \text{$p$ does not appear in $x$}\},\] 
    \[Y_p=\{y \in Y : \phi(y)\in X_p\}.\]
    
    Observe that $Y_p\subset Y$ and thus $\Lep_{W_k}(Y_p)\subset \Lep_{W_k}(Y)$. The number $k$ depends on $X$ and not on $p$, so we can ``hard-code'' it in the next algorithm, together with the finite set $\Lep_{W_k}(Y)$. 
    
    The key observation is that $p$ belongs to $\Lep(X)$ if and only if $\Lep_{W_k}(Y_p)$ is properly contained in $\Lep_{W_k}(Y)$. This is equivalent to the assertion ($p\not\in\Lep(X)$) $\iff$ $(\Lep_{W_k}(Y_p)=\Lep_{W_k}(Y))$. Let us prove this equivalence. 
    \begin{itemize}
        \item $(\Rightarrow)$. Suppose that $p\not\in\Lep(X)$. Then it follows from the definition of $X_p$ that $X_p=X$, which implies $Y_p=Y$, and therefore $\Lep_{W_k}(Y_p)=\Lep_{W_k}(Y)$. 
        \item $(\Leftarrow)$. Suppose that $\Lep_{W_k}(Y_p)=\Lep_{W_k}(Y)$. This assumption implies that we also have $\Le_{B_k}(Y_p)=\Le_{B_k}(Y)$.

        We first prove that $\phi(Y_p)=X$. Here we use that $X$ is projectively isolated, and our choice of $k$. The equality $\Le_{B_k}(Y_p)=\Le_{B_k}(Y)$ implies that $Y_p$ belongs to $U$. Then $\phi(Y_p)=X$ by definition of $U$. 
        
        Next we prove that $\phi(Y_p)=X_p$. Let $\psi\colon Y\to X$ be the restriction of $\phi$. Thus $\psi$ is a surjective function. It is clear that $\psi(Y_p)=\phi(Y_p)$, so it suffices to prove that $\psi(Y_p)=X_p$. We also have $\psi(\psi^{-1}(X_p))=X_p$ by an elementary property of surjective functions. But since $\psi^{-1}(X_p)=Y_p$, it follows that $\psi(Y_p)=X_p$ as claimed.   
        
        We proved that $\phi(Y_p)=X$ and also that $\phi(Y_p)=X_p$, so we can conclude that $X=X_p$. The last equality implies $p\not\in\Lep(X)$.
    \end{itemize} 
    We have proved our claim that $p\in\Lep(X)$ if and only if $\Lep_{W_k}(Y_p)\subsetneq\Lep_{W_k}(Y)$.
    
    In order to finish the proof, it suffices to show that the relation $\Lep_{W_k}(Y_p)\subsetneq\Lep_{W_k}(Y)$ is recursively enumerable. For this we will use the fact that $Y_p$ is effective uniformly on $p$, meaning that there is an algorithm which given $p\colon F\to \A$ ($F\Subset S^\ast$) and $q\colon E\to \B$ ($E\Subset S^\ast$) halts if and only if $q\in \Lep^c(Y_p)$. This will be proved separately in Lemma \ref{lemma} below. We will also use the fact that $\Lep_{W_k}(Y)$ is finite and does not depend on $p$. Let $\{q_1,\dots,q_m\}=\Lep_{W_k}(Y)$. In order to determine whether $\Lep_{W_k}(Y_p)\subsetneq\Lep_{W_k}(Y)$, it suffices to simultaneously run $m$ instances of the algorithm in the previous paragraph, one for each $q_i$, $i=1,\dots,m$. If some of these $m$ processes halts to conclude that $q_i\in\Lep^c_{W_k}(Y_p)$, then we stop the whole process to conclude $\Lep_{W_k}(Y_p)\subsetneq\Lep_{W_k}(Y)$ and consequently $p\in\Lep(X)$. \end{proof}

\end{document}
